\documentclass[10pt,a4paper]{amsart}
\usepackage[margin=19mm]{geometry}
\usepackage{amsmath,amssymb,mathtools}
\usepackage[scr=rsfs,cal=euler]{mathalfa}
\usepackage{xfrac}
\usepackage[unicode,hidelinks,bookmarks=false]{hyperref}
\newcommand{\ie}{i.e.\ }
\newcommand{\cf}{cf.\ }
\newcommand{\resp}{respectively\ }
\newcommand{\Subsection}{\subsection}
\providecommand{\nobreakdash}{\nobreak\hbox{-}}
\setlength{\emergencystretch}{2em}
\multlinegap0pt
\usepackage{amssymb}
\usepackage{xcolor}
\usepackage{enumerate}
\usepackage[shortlabels]{enumitem}
\usepackage{tikz}
\newtheorem{lem}{Lemma}
\newtheorem{prop}{Proposition}
\newtheorem{thm}{Theorem}
\newcommand{\hstack}{
  \mathbin{
    \tikz[baseline=-0.582ex]
      \draw[line width=.1ex, line cap=round]
        (0,0) circle (.75ex)
        (-.5ex,0) -- (.5ex,0);
  }
}
\newcommand{\subhstack}{
  \mathbin{
    \tikz[baseline=-0.449ex]
      \draw[line width=.075ex,line cap=round]
        (0,0) circle (.5625ex)
        (-.375ex,0)--(.375ex,0);
  }
}
\newcommand{\subvstack}{
  \mathbin{
    \tikz[baseline=-0.449ex]
      \draw[line width=.075ex,line cap=round]
        (0,0) circle (.5625ex)
        (0,-.375ex)--(0,.375ex);
  }
}
\newcommand{\vstack}{
  \mathbin{
    \tikz[baseline=-0.582ex]
      \draw[line width=.1ex, line cap=round]
        (0,0) circle (.75ex)
        (0,-.5ex) -- (0,.5ex);
  }
}
\title{A welded extension of \textit{4-}equivalence}
\author{Emmanuel Graff}
\date{Source: arXiv:2609.04624v2 (CC BY 4.0)}
\begin{document}
\maketitle

\noindent\emph{Source and licence.} A welded extension of \textit{4-}equivalence. Version arXiv:2609.04624v2 (CC BY 4.0), distributed by arXiv under the Creative Commons Attribution 4.0 International licence, http://creativecommons.org/licenses/by/4.0/, as recorded in the arXiv licence metadata for this exact version and shown on the abstract page https://arxiv.org/abs/2609.04624v2. Full licence notice: \texttt{licenses/arxiv-2609.04624-CC-BY-4.0.txt}.\par\smallskip

\noindent\textit{Editorial scope.} Counted unit: the article's thm thmtangleset (source0.tex lines 631--647). The article's complete original proof of it is delivered with it (source0.tex lines 649--695, one \texttt{proof} environment of 47 lines). No mathematics is written, repaired or re-derived: the statement and the proof are the article's own lines, in the article's own notation, and no retained line is substituted or re-flowed.  Retained uncounted as the source-local context the proof needs: lines 443--445; lines 459--461; lines 522--536; lines 604--614; lines 621--626.  One declared adaptation: exactly 3 ancillary strings inside the retained spans are omitted (image-inclusion commands and, where the source repeats a label, the repeated label definitions) and no other line differs from the source; the figure environments, with their placement, captions and labels, are kept, and the package records the omitted strings in fidelity receipts bound to the affected bodies.  No retained line contains a citation, so the wrapper carries no bibliography and no bibliography entry.  The retained lines contain the article's own diagram code, which the wrapper typesets with the standard tikz packages; no image file is delivered and the package declares no asset dependency.  Editorial formatting only, in addition: an AMS \texttt{amsart} preamble that restates the article's own declarations and macros used by the retained lines, namely the environments \texttt{lem}, \texttt{prop}, \texttt{thm} on separate counters, together with the standard packages those lines need.  The retained environments print in this document as Lemma 1, Proposition 1, Theorem 1 in order of appearance, whereas the article prints the same items with the numbering of its own full document.  The complete native source of the article as distributed in the arXiv e-print for this exact version is bundled unchanged as \texttt{sources/209377/source0.tex}.
\begin{lem}\label{lemarrowcommu}
Two parallel $w$-arrows with opposite orientations can be exchanged, up to wink-equivalence, by adding one twist to each of them.
\end{lem}

\begin{lem}\label{lemarrowsquareenclosed}
If the heads of two parallel $w$-arrows having the same orientation are enclosed by another $w$-arrow, then the enclosing $w$-arrow can be removed up to wink-equivalence.
\end{lem}

\begin{prop}\label{proptanglesets}
The sets $\mathcal{T}_{\subvstack}$ and $\mathcal{T}_{\subhstack}$ are finite and are given by
\[
\mathcal{T}_{\subvstack}
=
\{R^u \vstack L^v \vstack \rho^w \mid u,v\in\{0,1,2,3\},\; w\in\{0,1\}\},
\]
and
\[
\mathcal{T}_{\subhstack}
=
\{U^u \hstack D^v \hstack \rho^w \mid u,v\in\{0,1,2,3\},\; w\in\{0,1\}\}.
\]
The choice $u=v=w=0$ corresponds to the vertical $\infty$-tangle and the horizontal $0$-tangle, respectively, consisting of two parallel vertical strands and two parallel horizontal strands.
\end{prop}

\begin{prop}\label{proptantglesetintersection}
The intersection of $\mathcal{T}_{\subhstack}$ and $\mathcal{T}_{\subvstack}$ is
\[
\mathcal{T}_{\subhstack}\cap\mathcal{T}_{\subvstack}
=
\{T^u \hstack D^v \hstack \rho \mid u,v\in\{0,1,2,3\}\}
=
\{R^u \vstack L^v \vstack \rho \mid u,v\in\{0,1,2,3\}\}
.
\]
\end{prop}

\begin{figure}[!htbp]
\centering

\caption{Proof of Proposition~\ref{proptantglesetintersection}.}
\label{figproptantglesetintersection}
\end{figure}

\begin{thm}\label{thmtangleset}
The set $\mathcal{T}$ is finite and is given by
\[
\mathcal{T}
=
\{X(u,v),\,H(u,v),\,V(u,v)
\mid u,v\in\{0,1,2,3\}\},
\]
where
\[
X(u,v)=T^u \hstack D^v \hstack \rho,
\qquad
H(u,v)=T^u \hstack D^v,
\qquad
V(u,v)=R^u \vstack L^v.
\]
\end{thm}

\begin{proof}
The elementary tangles $\sigma$, $\overline{\sigma}$ and $\rho$ all belong to the set above. It is therefore enough to prove that this set is closed under the stacking operations $\vstack$ and $\hstack$, whenever no closed component is created.

By Proposition~\ref{proptantglesetintersection} and Lemma~\ref{lemarrowcommu}, rotation fixes the tangles $X(u,v)$ and exchanges $V(u,v)$ with $H(u,v)$. Since $\mathcal T$ is invariant under rotation, and rotation exchanges vertical and horizontal stacking, it is enough to establish closure under one of the two operations.

We therefore consider the nine possible horizontal stackings. The resulting configurations fall into four distinct cases, denoted by $A$, $B$, $C$, and $D$. Table~\ref{tablestack} lists the case corresponding to each of the nine stackings.

\begin{table}[!htbp]
\centering
\begin{tabular}{|c|c|c|c|}
\hline
$\hstack$ & $X(u,v)$ & $H(u,v)$ & $V(u,v)$\\
\hline
$X(u,v)$ & $A$ & $A$ & $C$\\
\hline
$H(u,v)$ & $A$ & $A$ & $B$\\
\hline
$V(u,v)$ & $C$ & $B$ & $D$\\
\hline
\end{tabular}
\vspace{0.3cm}
\caption{The nine possible horizontal stackings.}
\label{tablestack}
\end{table}
\vspace{-0.4cm}
The $D$-cases cannot occur, since such a product always creates a closed component. The four $A$-cases are covered by Proposition~\ref{proptanglesets}, because both factors belong to $\mathcal T_{\subhstack}$. The two $B$-cases are equivalent by rotation and follow essentially from the equalities in Figure~\ref{figthmproofBcases}.

\begin{figure}[!htbp]
\centering

\caption{Equalities used to prove the $B$-cases.}
\label{figthmproofBcases}
\end{figure}

The first equality follows from Tails Exchanges together with Isolated arrow when $u=0$, an arrow~wink-move when $u=\pm1$, and Lemma~\ref{lemarrowsquareenclosed} when $u=2$. The second is proved in the same way, with an additional application of Lemma~\ref{lemarrowcommu} to exchange the \emph{$u$-} and \emph{$v$-factors}.

The two $C$-cases are likewise equivalent by rotation. Their proof combines the equalities of Figure~\ref{figthmproofBcases} with those of Figure~\ref{figthmproofCcases}.

\begin{figure}[!htbp]
\centering

\caption{Equalities used to prove the $C$-cases.}
\label{figthmproofCcases}
\end{figure}

The first equality follows from the Arrow isotopy, Head/Tail Reversal, and Tails Exchange moves. A similar version of the second equality appears in Figure~\ref{figproptantglesetintersection} and was discussed in the proof of Proposition~\ref{proptantglesetintersection}.
\end{proof}

\end{document}
